\section{EXPERIMENTAL METHODOLOGY}

All primary results must be obtained under onboard battery power unless an
experiment explicitly studies tethered operation. Before testing, the final mass,
battery state, tire condition, firmware, controller revision, and surface must be
recorded. Each quantitative condition should use at least ten independent trials
or justify a smaller number.

\begin{table*}[t]
  \caption{Planned experiments and minimum reported quantities}
  \label{tab:protocol}
  \centering
  \small
  \begin{tabular}{p{0.18\textwidth}p{0.24\textwidth}p{0.28\textwidth}p{0.20\textwidth}}
    \toprule
    Experiment & Controlled condition & Primary metrics & Required evidence \\
    \midrule
    Untethered operation & fixed battery state and payload & runtime, voltage sag, failures & continuous raw video and log \\
    Flat-ground walking & measured course, common surface & speed, success, energy/distance & timestamps and distance scale \\
    Flat-ground rolling & same course and payload & speed, success, energy/distance & same protocol as walking \\
    Single stair & measured riser and tread & time, success, pitch, peak current & repeated trials, failures included \\
    Continuous staircase & 15 measured steps & ascent time, success, pitch/roll & uncut full-ascent video \\
    Irregular terrain & defined obstacle field & success, speed, failure mode & terrain dimensions and trial list \\
    Mode transition & robot initially stationary & transition time, success, energy & uncut transition sequence \\
    \bottomrule
  \end{tabular}
\end{table*}

\subsection{Same-Hardware Ablations}

Three controller restrictions separate the effect of the morphology from changes
between prototype generations: (A) distal-only rotating-leg motion with the
articulated joints locked, (B) articulated walking with distal rotation disabled,
and (C) the complete proposed controller. Conditions A--C use the same robot,
battery, tires, payload, and environment.

\subsection{Metrics and Statistics}

For a course of measured length $d$ completed in time $t$, speed is $v=d/t$.
Success rate over $N$ trials is $p_s=N_{\mathrm{success}}/N$. Report mean,
standard deviation, all $N$, and individual values when $N$ is small. Electrical
energy is computed from synchronized voltage and current measurements,
\begin{equation}
  E = \int_{t_0}^{t_1} V(t)I(t)\,dt,
\end{equation}
and normalized by distance or gained height as appropriate.

\draftnote{Pre-register exact start/finish events for every timing result. Resolve
all historical v1 speed inconsistencies before using v1 as supporting evidence.}
